\documentclass[12pt, a4paper]{article}
\usepackage[utf8]{inputenc}
\usepackage[T1]{fontenc}
\usepackage[romanian]{babel}
\usepackage{amsmath}
\usepackage{amssymb}
\usepackage{geometry}
\geometry{a4paper, margin=2cm}
\usepackage{listings}
\usepackage{xcolor}

\lstset{
    language=Python,
    basicstyle=\ttfamily\small,
    keywordstyle=\color{blue},
    stringstyle=\color{red},
    commentstyle=\color{green!50!black},
    showstringspaces=false,
    breaklines=true,
    frame=single,
    numbers=left,
    numberstyle=\tiny\color{gray}
}

\title{Rezolvare Teme - Algoritmi de Primalitate \c{s}i Factorizare}
\author{}
\date{}

\begin{document}
\maketitle
\section*{Tema 6 (Subpunctul 43): Factorizarea num\u{a}rului $7373$}

\textbf{Cerin\c{t}\u{a}}: Factorizarea lui $n = 7373$ folosind metoda Fermat.

\begin{enumerate}
    \item Calcul\u{a}m $t = \lceil\sqrt{n}\rceil$: \\
    $\sqrt{7373} \approx 85.86 \implies t_{start} = 86$.
    \item Verific\u{a}m dac\u{a} $t^2 - n$ este p\u{a}trat perfect ($s^2$):
    \begin{itemize}
        \item Pentru $t = 86$: \\
        $86^2 - 7373 = 7396 - 7373 = 23$ (Nu este p\u{a}trat perfect).
        \item Pentru $t = 87$: \\
        $87^2 - 7373 = 7569 - 7373 = 196$.
    \end{itemize}
    \item Deoarece $196 = 14^2$, am g\u{a}sit $s = 14$.
    \item Calcul\u{a}m factorii $a$ \c{s}i $b$ folosind $t = 87$ \c{s}i $s = 14$:
    \begin{align*}
        a &= t - s = 87 - 14 = 73 \\
        b &= t + s = 87 + 14 = 101
    \end{align*}
    \end{enumerate}

\textbf{Rezultat final}: $7373 = 73 \times 101$.


\end{document}